\section{Continuous-magnitude observed-state lower bounds}
\label{app:continuous-observed}

This appendix verifies the lower bounds and regularity assertions of
\cref{prop:continuous}. The class permits arbitrary finite measurable
sign-cell partitions and disclosed gap maps. The query policy may choose
any state, revisit states, and stop after an unbounded number of observations.
The upper-bound transfer from fresh outer sampling remains a paper proof.

\begin{theorem}[Verified continuous-magnitude obstruction]
\label{thm:continuous-observed-verified}
Let $U$ be uniform on $[0,1]$ and let $X\in\{-1,1\}$ have a deterministic
conditional Bernoulli law with
$|\mathbb E[X\mid U]|=2D(U)$ and
$D(U)\sim\operatorname{Unif}(0,1/4]$.
There is a class of such models with arbitrary finite sign-cell
partitions for which, for all sufficiently small $\varepsilon>0$,
the minimax expected observation counts satisfy
\[
\begin{gathered}
 \mathcal C_{\rm conf}(\varepsilon,1/16)
 \ge c\,\varepsilon^{-2}\log^2(1/\varepsilon),
 \qquad
 \mathcal C_{\rm RMS}(\varepsilon)
 \ge \frac{c}{64}\,\varepsilon^{-2}\log^2(1/\varepsilon),\\
 c=\frac{225\,\operatorname{kl}(7/8,1/8)}{128\cdot64^2}>0.
\end{gathered}
\]
Every model has a version with $0<D(u)\le1/4$ at every state.
Its normalized conditional magnitude has a density bounded by $2$.
For every real $q\ge0$ and every state, its normalized central
$q$th absolute moment is at most $3^{q/2}$.
\end{theorem}

\begin{proof}
Use the finite measurable design from \cref{app:finite-observed},
retaining the unrounded gap coordinate as the actual magnitude.
The sign-cell caps still determine the testing priors, while the
actual gap law is exactly the prescribed law $\mu$.
For $\mu=\operatorname{Unif}(0,1/4]$, cap both the gap and each cell
magnitude at $1/4$. This preserves the gap law, increases the clipped
modulus of the caps, and retains the target-concentration condition.
It also enforces the cutoff at states of probability zero that a
selected-query algorithm could otherwise choose.

For a point with actual gap $d$ in a cell with cap $b\ge d$, obtain
one observation from the $b$-gap cell coin and independently flip its
outcome with probability
\[
 \gamma(b,d)=\frac{1-d/b}{2}.
\]
For either cell sign, the resulting Bernoulli parameter is exactly
$1/2\pm d$. A history decoder chooses each subsequent point from the
transformed observations. It preserves every paid-observation
indicator and the total work on every tape.

The checked simulation extracts one raw cell observation and one
independent uniform noise value at each active step. The transformed
observation has the requested conditional law, and the unused
randomness retains its independent product law. Induction therefore
gives the ordinary fresh-Bernoulli law for every finite adaptive
prefix, jointly with the original private seed. These prefix laws
uniquely determine the complete transcript measure.
Consequently the finite-cell testing argument applies to every
physical experiment with this observation rule, giving
\[
 B\ge\frac{225}{128h^2}\operatorname{kl}(1-2\delta,2\delta),
 \qquad \Delta_h\ge\frac18\omega_\mu(h).
\]

For $0<h<1/4$, the exact uniform-law modulus is
\[
 \omega_\mu(h)=4h\left(1+\log\frac1{4h}\right).
\]
For sufficiently small $\varepsilon$, the checked choice
$h=64\varepsilon/\log(1/\varepsilon)$ has $h\le1/64$ and
$\varepsilon<\omega_\mu(h)/32\le\Delta_h/4$.
Take $\delta=1/16$ to obtain the stated confidence coefficient.
A squared-error guarantee $\varepsilon^2$ gives radius $4\varepsilon$
at failure probability $1/16$ by Markov's inequality.
Apply the confidence result at $4\varepsilon$ and use
$\log(1/(4\varepsilon))\ge\frac12\log(1/\varepsilon)$
for $\varepsilon\le1/16$ to obtain the RMS coefficient.
\end{proof}

The regularity calculation concerns the actual conditional law.
Writing $m=\mathbb E[X\mid U]$, direct Bernoulli integration gives
\[
 |m|=2D,\qquad
 \operatorname{Var}(X\mid U)=1-4D^2\ge\frac34,\qquad
 |X-m|\le\frac32.
\]
Hence $|X-m|/\sqrt{\operatorname{Var}(X\mid U)}\le\sqrt3$,
which proves the moment assertion.
For the normalized magnitude, put
\[
 Z(d)=\frac{2d}{\sqrt{1-4d^2}},\qquad
 G(z)=\frac{z}{2\sqrt{1+z^2}},\qquad
 G'(z)=\frac1{2(1+z^2)\sqrt{1+z^2}}.
\]
The inverse $G$ is strictly increasing and $0<G'\le1/2$.
On $0<d\le1/4$, $G(Z(d))=d$. The checked change-of-variables
identity thus gives the density
\[
 f_Z(z)=
 \frac{2\,\mathbf1_{(0,Z(1/4)]}(z)}
 {(1+z^2)\sqrt{1+z^2}}
 \le2.
\]
These uniform bounds verify the bounded-density and normalized
conditional-moment assumptions stated in \cref{prop:continuous}.

The \href{https://github.com/davidiach/adaptive-threshold-estimation/blob/10fc9a912ff5ecd9b1125bea47daa4821df6fd79/lean/ThresholdHC/ContinuousObserved/UniformLower.lean}{confidence theorem},
\href{https://github.com/davidiach/adaptive-threshold-estimation/blob/10fc9a912ff5ecd9b1125bea47daa4821df6fd79/lean/ThresholdHC/ContinuousObserved/UniformRMSLower.lean}{RMS theorem}, and
\href{https://github.com/davidiach/adaptive-threshold-estimation/blob/10fc9a912ff5ecd9b1125bea47daa4821df6fd79/lean/ThresholdHC/ContinuousObserved/Regularity.lean}{regularity theorem}
identify the physical observation law, threshold target, gap
distribution, normalized density, and conditional moments.
The finite replication count uses the second-moment construction
from \cref{app:finite-observed}. The class contains designs at all
scales and is selected before the requested accuracy.
